\documentclass[conference]{IEEEtran}
\IEEEoverridecommandlockouts

\usepackage{cite}
\usepackage{amsmath,amssymb,amsfonts}
\usepackage{algorithmic}
\usepackage{graphicx}
\usepackage{textcomp}
\usepackage{xcolor}
\usepackage{booktabs}
\usepackage{multirow}
\usepackage{url}
\usepackage{hyperref}

\begin{document}

% =========================================================
% TITLE AND AUTHORS
% =========================================================

\title{Your Proposed Research Title}

\author{
\IEEEauthorblockN{Student Name 1, Student Name 2, Student Name 3}
\IEEEauthorblockA{
Department of Computer Science and Engineering\\
Mississippi State University\\
Starkville, MS, USA\\
\{email1, email2, email3\}@msstate.edu
}
{\color{blue} For Phase 2, update all black-colored content in this template within your existing Overleaf project.}}

\maketitle



% =========================================================
% ABSTRACT
% =========================================================

\begin{abstract}

% Address the following in approximately 150--200 words:
%
% 1. CONTEXT:
%    What general research area does this work address?
%
% 2. PROBLEM:
%    What specific problem is being studied?
%
% 3. GAP:
%    What important limitation remains in existing work?
%
% 4. PROPOSED APPROACH:
%    What method/model/framework do you propose?
%
% 5. EVALUATION:
%    What dataset(s), baselines, and metrics will be used?
%
% 6. EXPECTED/FINAL RESULT:
%    During the proposal stage, state what you intend to evaluate.
%    In the final paper, replace this with quantitative results.
%
% 7. CONTRIBUTION:
%    What is the main contribution of the study?
%
% NOTE:
% Do not make unsupported claims such as
% "our method significantly outperforms existing methods"
% until experimental evidence is available.

[Write your abstract here.]

\end{abstract}


\begin{IEEEkeywords}
machine learning, [keyword 2], [keyword 3], [keyword 4]
\end{IEEEkeywords}


% =========================================================
% I. INTRODUCTION
% =========================================================

\section{Introduction}

% ---------------------------------------------------------
% Paragraph 1: Background and Context
% ---------------------------------------------------------
%
% Address:
% - What is the broader application/research area?
% - What ML problem is involved?
% - Why does this problem matter?
% - Who or what benefits from solving it?
% - Support important claims with references.

[Introduce the research area and application context.]


% ---------------------------------------------------------
% Paragraph 2: Importance of the Problem
% ---------------------------------------------------------
%
% Address:
% - What is the specific problem?
% - Why is solving it important?
% - What happens if the problem is not addressed?
% - Can its importance be quantified?
%
% Examples of useful evidence:
% - error rates
% - class imbalance
% - computational cost
% - robustness issues
% - latency
% - fairness
% - generalization
% - data scarcity
% - interpretability

[Explain the specific problem and its importance.]


% ---------------------------------------------------------
% Paragraph 3: Existing Work
% ---------------------------------------------------------
%
% Address:
% - What approaches have already been used?
% - What are the strongest/relevant existing methods?
% - What have previous studies achieved?
%
% Cite the most relevant papers.

[Briefly summarize existing approaches.]


% ---------------------------------------------------------
% Paragraph 4: Research Gap
% ---------------------------------------------------------
%
% Clearly identify the gap.
%
% Use evidence from the literature rather than simply stating:
% "No one has done this before."
%
% A strong gap statement should establish:
%
% Existing literature -> limitation -> consequence -> need
%
% Example structure:
% Although existing methods achieve ______,
% they remain limited by ______.
% This limitation becomes important when ______.
% Therefore, there is a need for ______.

[State the research gap.]


% ---------------------------------------------------------
% Paragraph 5: Proposed Solution
% ---------------------------------------------------------
%
% Address:
% - What do you propose?
% - Which gap/challenge does it address?
% - Why should the proposed approach help?

[Introduce the proposed approach.]


% ---------------------------------------------------------
% Contributions
% ---------------------------------------------------------

The main contributions of this work are summarized as follows:

\begin{itemize}

    \item \textbf{Contribution 1:}
    [State the first specific contribution.]

    \item \textbf{Contribution 2:}
    [State the second specific contribution.]

    \item \textbf{Contribution 3:}
    [State the third specific contribution.]

\end{itemize}

% IMPORTANT:
% Contributions should describe what YOUR study contributes,
% not general statements such as:
% "We study machine learning."
%
% Each contribution should eventually be supported by
% experiments, analysis, or theoretical results.


% =========================================================
% II. RELATED WORK
% =========================================================

\section{Related Work}

% Organize related work by METHODS/CHALLENGES rather than
% writing one paragraph per paper.
%
% Possible structure:
%
% A. Traditional ML Approaches
% B. Deep Learning Approaches
% C. Recent/State-of-the-Art Approaches
%
% For every group of papers discuss:
%
% 1. What problem do they address?
% 2. What dataset do they use?
% 3. What method do they use?
% 4. What evaluation metrics are reported?
% 5. What are their important findings?
% 6. What limitations remain?
%
% The objective is NOT only to summarize papers.
% The objective is to establish the research gap.

\subsection{Existing Approaches}

[Discuss the major categories of existing methods.]


\subsection{Limitations of Existing Work}

[Compare the limitations that motivate your proposed study.]


% ---------------------------------------------------------
% Literature Comparison Table
% ---------------------------------------------------------

\begin{table*}[t]
\centering
\caption{Comparison of Representative Related Works}
\label{tab:literature}
\resizebox{\textwidth}{!}{
\begin{tabular}{lllllll}
\toprule
\textbf{Reference} &
\textbf{Dataset} &
\textbf{Method} &
\textbf{Task} &
\textbf{Metrics} &
\textbf{Key Finding} &
\textbf{Limitation / Gap} \\
\midrule

Paper 1 & Dataset & Method & Task & Metric & Finding & Gap \\

Paper 2 & Dataset & Method & Task & Metric & Finding & Gap \\

Paper 3 & Dataset & Method & Task & Metric & Finding & Gap \\

Paper 4 & Dataset & Method & Task & Metric & Finding & Gap \\

Paper 5 & Dataset & Method & Task & Metric & Finding & Gap \\
.... & .... & .... & .... & .... & .... \\
\bottomrule
\end{tabular}}
\end{table*}


% =========================================================
% III. RESEARCH CHALLENGES AND MOTIVATION
% =========================================================

\section{Research Challenges and Motivation}

% Identify THREE important challenges derived from your
% literature survey.
%
% Each challenge MUST contain:
%
% 1. What is the challenge?
% 2. Evidence that the challenge exists.
% 3. Why is it important?
% 4. Why is it technically difficult?
% 5. Which existing methods are affected?
% 6. What happens if it is not addressed?
% 7. Which proposed component will address it?


\subsection{Challenge 1: [Challenge Name]}

\textbf{Problem:}
[Precisely define the challenge.]

\textbf{Evidence from literature:}
[Cite papers/results demonstrating the problem.]

\textbf{Importance:}
[Explain why addressing this challenge matters.]

\textbf{Technical difficulty:}
[Explain specifically why this challenge is difficult.]

\textbf{Connection to proposed method:}
[Identify the proposed component intended to address it.]


\subsection{Challenge 2: [Challenge Name]}

\textbf{Problem:}
[Define Challenge 2.]

\textbf{Evidence from literature:}
[Provide supporting evidence and citations.]

\textbf{Importance:}
[Explain importance.]

\textbf{Technical difficulty:}
[Explain technical difficulty.]

\textbf{Connection to proposed method:}
[Identify the corresponding proposed technique.]


\subsection{Challenge 3: [Challenge Name]}

\textbf{Problem:}
[Define Challenge 3.]

\textbf{Evidence from literature:}
[Provide supporting evidence and citations.]

\textbf{Importance:}
[Explain importance.]

\textbf{Technical difficulty:}
[Explain technical difficulty.]

\textbf{Connection to proposed method:}
[Identify the corresponding proposed technique.]


% =========================================================
% IV. PROPOSED METHODOLOGY
% =========================================================

\section{Proposed Methodology}

% Start this section with the overall research idea.
%
% Answer:
% - What is being proposed?
% - What is the input?
% - What is the output?
% - What are the major components?
% - How do the components interact?

%%% More details 
% Explain the complete method from INPUT to OUTPUT.
%
% Address:
%
% 1. What data enters the system?
% 2. What preprocessing is performed?
% 3. What features/representations are used?
% 4. What is the baseline architecture?
% 5. What new/modified component is proposed?
% 6. Where is the proposed component inserted?
% 7. What does each component mathematically/functionally do?
% 8. Which research challenge does each component address?
% 9. What loss/objective function is optimized?
% 10. What output does the model produce?
% 11. How is inference/prediction performed?
% 12. Why was this design chosen instead of alternatives?
%



\subsection{Problem Formulation}

% Clearly define:
%
% Input:
% X = ?
%
% Output/target:
% y = ?
%
% Dataset:
% D = {(x_i,y_i)}?
%
% Learning objective:
% classification / regression / clustering / forecasting /
% representation learning / etc.
%
% Define mathematical notation when appropriate.

[Formulate the ML problem.]


\subsection{Overall Framework}

% INSERT A SYSTEM/METHOD DIAGRAM.
%
% Recommended flow:
%
% Raw Data
%    ->
% Preprocessing
%    ->
% Feature Representation
%    ->
% Baseline/Proposed Model
%    ->
% Prediction
%    ->
% Evaluation
%
% Every major block in the figure must be explained in text.

\begin{figure}[htbt]
    \centering
    % \includegraphics[width=\columnwidth]{figures/framework.pdf}
    \fbox{
        \parbox[c][3cm][c]{0.9\columnwidth}{
        \centering
        Insert proposed methodology / ML pipeline diagram here
        }
    }
    \caption{Overview of the proposed methodology.}
    \label{fig:framework}
\end{figure}


\subsection{Proposed Components}

% Explain each component.
%
% For EACH proposed component answer:
%
% 1. What does it do?
% 2. What input does it receive?
% 3. What output does it produce?
% 4. Which research challenge does it address?
% 5. Why was this technique selected?
% 6. How is it different from the baseline/existing approach?

[Describe the proposed components.]


% ---------------------------------------------------------
% Challenge-to-Method Mapping
% ---------------------------------------------------------

\begin{table}[t]
\centering
\caption{Mapping Between Research Challenges and Proposed Methods}
\label{tab:challenge_method}
\resizebox{\columnwidth}{!}{
\begin{tabular}{llll}
\toprule
\textbf{Challenge} &
\textbf{Technique} &
\textbf{Function} &
\textbf{Expected Effect} \\
\midrule

Challenge 1 & Technique 1 & [Function] & [Effect] \\

Challenge 2 & Technique 2 & [Function] & [Effect] \\

Challenge 3 & Technique 3 & [Function] & [Effect] \\

\bottomrule
\end{tabular}}
\end{table}


% =========================================================
% V. EXPERIMENTAL METHODOLOGY
% =========================================================





% =========================================================
% VII. SIX-WEEK EXECUTION PLAN
% =========================================================

\section{Project Execution Plan}


%=========================================================
% VII. SIX-WEEK RESEARCH AND IMPLEMENTATION PLAN
% =========================================================

\subsection{Six-Week Research and Implementation Plan}

% Provide a realistic six-week plan for completing the proposed study.
%
% The plan should demonstrate how the project progresses from
% literature analysis and baseline implementation to the proposed
% method, experiments, analysis, and final manuscript.
%
% For EACH week, specify:
%
% 1. Research/technical objective
% 2. Specific tasks to be completed
% 3. Expected measurable output/deliverable
% 4. Team member(s) responsible
%
% Avoid vague descriptions such as:
% "Work on model," "Do experiments," or "Write paper."
%
% Tasks should be specific and verifiable.
%
% Example:
% Weak: "Implement models."
%
% Better:
% "Implement logistic regression and random forest baselines,
% evaluate using stratified 5-fold cross-validation, and report
% Accuracy, Precision, Recall, and F1-score."


\begin{table*}[h!]
\centering
\caption{Six-Week Research and Implementation Plan Example}
\label{tab:six_week_plan}
\renewcommand{\arraystretch}{1.25}

\begin{tabular}{
p{0.05\textwidth}
p{0.16\textwidth}
p{0.31\textwidth}
p{0.28\textwidth}
p{0.12\textwidth}
}
\toprule

\textbf{Week} &
\textbf{Objective} &
\textbf{Specific Tasks} &
\textbf{Expected Output / Milestone} &
\textbf{Member(s)} \\

\midrule

1 &
Problem and data preparation &
Finalize research questions and hypotheses; finalize dataset;
implement preprocessing; define train/validation/test protocol;
identify evaluation metrics. &
Reproducible dataset pipeline; finalized research questions;
dataset statistics and experimental protocol. &
[Names] \\

\midrule

2 &
Baseline establishment &
Implement selected baseline methods; verify implementations;
perform initial training and evaluation using the same
experimental protocol planned for the proposed method. &
Baseline performance table; initial plots; verified baseline
implementation. &
[Names] \\

\midrule

3 &
Proposed method implementation &
Implement the proposed technique/component(s);
integrate them with the existing data pipeline;
perform debugging and preliminary experiments. &
Working implementation of proposed method;
preliminary results demonstrating successful execution. &
[Names] \\

\midrule

4 &
Main experiments &
Train and evaluate the proposed method;
perform hyperparameter selection;
compare against baselines using the predefined metrics;
repeat experiments where necessary. &
Main experimental results; comparison tables;
performance figures. &
[Names] \\

\midrule

5 &
Analysis and validation &
Conduct ablation study, error/failure analysis,
robustness or sensitivity analysis where appropriate;
analyze why the proposed method improves or fails. &
Ablation/error analysis; final figures and tables;
evidence supporting or rejecting the research hypothesis. &
[Names] \\

\midrule

6 &
Finalization and reproducibility &
Complete remaining experiments; verify reproducibility;
interpret results; identify limitations;
finalize manuscript, figures, references, and presentation. &
Reproducible final results; completed IEEE-style manuscript;
final presentation/poster. &
[Names] \\

\bottomrule
\end{tabular}

\end{table*}
% =========================================================
% VIII. TEAM MEMBER CONTRIBUTIONS
% =========================================================

\section{Team Member Contributions}

% Clearly state each person's responsibilities.
%
% Avoid:
% "Everyone worked equally."
%
% Instead identify concrete contributions:
%
% - literature survey
% - dataset preparation
% - preprocessing
% - baseline implementation
% - proposed algorithm
% - experiments
% - statistical analysis
% - visualization
% - manuscript preparation

\begin{itemize}

    \item \textbf{Student 1:}
    [Specific contributions.]

    \item \textbf{Student 2:}
    [Specific contributions.]

    \item \textbf{Student 3:}
    [Specific contributions.]

\end{itemize}


\newpage

\color{red}
\section{Experimental Methodology}


\subsection{Dataset}

% Provide:
%
% Dataset name:
% Dataset source/URL:
% Public/private:
% Number of samples:
% Number of features:
% Number of classes/targets:
% Class distribution:
% Missing values:
% Data modality:
%     tabular/image/text/audio/time-series/etc.
%
% Explain why this dataset is suitable for the research question.

[Describe the dataset.]


\subsection{Data Preprocessing}

% Explain ALL relevant preprocessing:
%
% - missing-value handling
% - normalization/standardization
% - categorical encoding
% - resizing
% - tokenization
% - augmentation
% - feature extraction
% - outlier treatment
% - class balancing
%
% Avoid data leakage.
%
% Clearly state which operations are learned only from
% the training set.

[Describe preprocessing.]


\subsection{Data Splitting and Validation}

% Specify:
%
% Training set: ___%
% Validation set: ___%
% Test set: ___%
%
% OR
%
% K-fold cross-validation with K = ___
%
% Explain WHY this evaluation strategy is appropriate.
%
% For imbalanced classification, consider stratification.
% For time-series data, do not randomly mix future samples
% into the training set.

[Describe the validation protocol.]


\subsection{Baseline Methods}

% Include meaningful comparison methods.
%
% Recommended:
%
% Baseline 1: simple/traditional baseline
% Baseline 2: stronger ML/DL baseline
% Baseline 3: relevant recent method
%
% Explain why each baseline is included.

\begin{itemize}

    \item \textbf{Baseline 1:} [Method and justification.]

    \item \textbf{Baseline 2:} [Method and justification.]

    \item \textbf{Baseline 3:} [Method and justification.]

\end{itemize}


\subsection{Implementation Details}

% Report enough information to make the work reproducible:
%
% Programming language:
% ML framework:
% Hardware:
% Optimizer:
% Learning rate:
% Batch size:
% Number of epochs:
% Loss function:
% Random seed(s):
% Hyperparameter search:
% Early stopping:
%
% Include only parameters relevant to your study.

[Provide implementation details.]


\subsection{Evaluation Metrics}

% List the PRIMARY and SECONDARY evaluation metrics.
%
% Classification examples:
% Accuracy, Precision, Recall, F1-score, AUROC
%
% Regression examples:
% MAE, MSE, RMSE, R^2
%
% Clustering examples:
% Silhouette score, ARI, NMI
%
% Explain WHY the selected metrics are appropriate.
%
% Do not report accuracy alone for a highly imbalanced
% classification problem without justification.

[Define and justify evaluation metrics.]


% =========================================================
% VI. RESULTS AND ANALYSIS
% =========================================================

\section{Results and Analysis}

% PROPOSAL STAGE:
% Leave the numerical entries blank or write "To be evaluated."
%
% FINAL PAPER:
% Replace with actual experimental results.


\subsection{Main Results}

\begin{table}[t]
\centering \color{red}
\caption{Performance Comparison}
\label{tab:results}
\begin{tabular}{lccc}
\toprule
\textbf{Method} &
\textbf{Metric 1} &
\textbf{Metric 2} &
\textbf{Metric 3} \\
\midrule

Baseline 1 & -- & -- & -- \\
Baseline 2 & -- & -- & -- \\
Baseline 3 & -- & -- & -- \\
Proposed Method & -- & -- & -- \\

\bottomrule
\end{tabular}
\end{table}


% For each major result discuss:
%
% 1. What was observed?
% 2. How large is the improvement/degradation?
% 3. Why might this behavior occur?
% 4. Does it support the original hypothesis?
% 5. How does it compare with the baseline?
%
% Do not only repeat numbers from the table.

[Analyze the main results.]


\subsection{Ablation Study}

% If your proposed approach contains multiple components,
% evaluate their individual contributions.
%
% Example:
%
% Baseline
% Baseline + Component A
% Baseline + Component B
% Full proposed method

\begin{table}[t]
\centering \color{red}
\caption{Ablation Study}
\label{tab:ablation}
\begin{tabular}{lc}
\toprule
\textbf{Configuration} & \textbf{Primary Metric} \\
\midrule

Baseline & -- \\
+ Component A & -- \\
+ Component B & -- \\
Full Proposed Method & -- \\

\bottomrule
\end{tabular}
\end{table}


\subsection{Error and Failure Analysis}

% Investigate WHERE the model fails.
%
% Possible analyses:
%
% - confusion matrix
% - difficult classes
% - false positives/false negatives
% - performance across subgroups
% - sensitivity to noise
% - sensitivity to training size
% - robustness
% - outliers
%
% Explain what the failures reveal about the model.

[Provide error/failure analysis.]



% =========================================================
% IX. USE OF GENERATIVE AI / LLMs
% =========================================================


\section{Use of Generative AI and LLM Tools}

% Disclose any use of ChatGPT, Copilot, Gemini,
% Claude, or other generative-AI/LLM tools.
%
% Address:
%
% 1. Tool/model used:
%
% 2. What was it used for?
%    Examples:
%    - brainstorming
%    - literature organization
%    - code assistance
%    - debugging
%    - writing assistance
%    - LaTeX formatting
%
% 3. Which parts of the work were AI-assisted?
%
% 4. How were generated outputs verified?
%
% 5. Were references/citations independently verified?
%
% 6. Were experimental results generated by your actual
%    implementation rather than accepted from an LLM?
%
% Students remain responsible for the correctness of
% submitted content.


[Provide your disclosure here.]


% =========================================================
% X. CONCLUSION
% =========================================================

\section{Conclusion}

% PROPOSAL:
% Summarize:
% - problem
% - identified gap
% - proposed solution
% - planned evaluation
% - expected contribution
%
% FINAL PAPER:
% Replace expected outcomes with actual findings,
% limitations, and possible future work.

[Write the conclusion.]


% =========================================================
% REFERENCES
% =========================================================

\bibliographystyle{IEEEtran}
\bibliography{references}

\end{document}